\documentclass[conference]{IEEEtran}
\usepackage{times}
\usepackage[numbers]{natbib}
\usepackage{graphicx}
\usepackage{booktabs}
\usepackage[bookmarks=true]{hyperref}
\usepackage{tabularx}
\begin{document}
\small

\title{Rebuttal for CoRL Submission \#416\vspace{-0.9cm}}
\maketitle
\thispagestyle{empty}
\vspace{-0.15cm}

{\small We thank the AC and all reviewers for constructive comments.
We clarify review questions as follows (R1=BKKY, R2=dWbG, R3=4zfY). }

\section*{1. Additional Baselines \& Downstream Value}

\textbf{Q1 (AC, R2, R3): Do we compare with ActNeRF/PB-NBV, and how does AURORA relate to Scalable Real2Sim?}
We have compared AURORA with adapted ActNeRF~\cite{actnerf} and
PB-NBV~\cite{pbnbv}; performance is reported in Table~\ref{tab:planner}.
Ray-GPIS plans faster and ranks executable actions better because it scores
reconstruction uncertainty directly under the in-hand action constraints. We
do not benchmark Scalable Real2Sim~\cite{pfaff2025_scalable_real2sim} because
it uses prescribed scanning trajectories rather than adaptive view planning;
its acquisition strategy is closer to the pose-conditioned baseline in Q2.

% =========================================================
% Table 1: Planner comparison
% =========================================================
\begin{table}[t]
\centering
\caption{Planner comparison over 120 paired episodes with GT pose and exact
RGB--depth registration.}
\label{tab:planner}
\vspace{-0.12cm}

\footnotesize
\setlength{\tabcolsep}{3pt}
\renewcommand{\arraystretch}{0.96}

\begin{tabular}{lcc}
\toprule
\multicolumn{3}{l}{\textbf{Reconstruction performance}} \\
Method & F-AUC$\uparrow$ & Final F$\uparrow$ \\
\midrule
Fixed
  & $0.7948{\pm}0.0061$ & $0.9401{\pm}0.0046$ \\
ActNeRF~\cite{actnerf}
  & $0.8665{\pm}0.0062$ & $0.9651{\pm}0.0050$ \\
PB-NBV~\cite{pbnbv}
  & $0.7475{\pm}0.0193$ & $0.8496{\pm}0.0234$ \\
Pose-Novelty
  & $0.9090{\pm}0.0043$ & $0.9945{\pm}0.0014$ \\
\textbf{Ray-GPIS}
  & $\mathbf{0.8987{\pm}0.0091}$ & $\mathbf{0.9814{\pm}0.0100}$ \\
\bottomrule
\end{tabular}

\vspace{3pt}

\begin{tabular}{lccc}
\toprule
\multicolumn{4}{l}{\textbf{Planner ranking and efficiency}} \\
Method & Corr.$\uparrow$ & Regret$\downarrow$ & Time (s)$\downarrow$ \\
\midrule
Fixed
  & -- & $0.0806{\pm}0.0025$ & $0.000$ \\
ActNeRF
  & $0.1619{\pm}0.0337$
  & $0.0317{\pm}0.0029$
  & $2.238{\pm}0.018$ \\
PB-NBV
  & $-0.3563{\pm}0.0731$
  & $0.1076{\pm}0.0124$
  & $0.391{\pm}0.024$ \\
Pose-Novelty
  & $0.6186{\pm}0.0294$
  & $0.0025{\pm}0.0006$
  & $2.49{\times}10^{-4}$ \\
\textbf{Ray-GPIS}
  & $\mathbf{0.4722{\pm}0.0667}$
  & $\mathbf{0.0088{\pm}0.0051}$
  & $\mathbf{0.261{\pm}0.024}$ \\
\bottomrule
\end{tabular}

\vspace{-0.12cm}
\end{table}

\textbf{Q2 (AC, R1, R2): Would a simpler pose-conditioned reorientation
achieve similar performance?} Yes. With exact poses and RGB--depth registration,
Pose-Novelty slightly outperforms AURORA (Table~\ref{tab:planner}); when a
visited view's depth is not fused, it underperforms AURORA
(Table~\ref{tab:robust}(a)) because it marks the pose as covered while AURORA
plans from the geometry actually acquired. AURORA's sensitivity to 6D pose
errors is summarized in Q5.

% =========================================================
% Table 2: Robustness and ablations
% =========================================================
\begin{table}[t]
\centering
\caption{Robustness and ablation results: \textbf{(a) Baseline robustness}
under sensing failure; \textbf{(b) Ablation} of Ray-GPIS components.}
\label{tab:robust}
\vspace{-0.12cm}

\footnotesize
\setlength{\tabcolsep}{3pt}
\renewcommand{\arraystretch}{0.96}

\begin{tabular}{p{0.44\columnwidth}lcc}
\toprule
Test / comparator & Metric & Comp. & \textbf{Full} \\
\midrule

\multicolumn{4}{l}{\textbf{(a) Baseline robustness}} \\
Visited but unfused (70\%) / Pose-Novelty
  & R@5$\uparrow$
  & 0.2827
  & $\mathbf{0.5066}$ \\

\addlinespace[3pt]
\multicolumn{4}{l}{\textbf{(b) Ablation}} \\
Clean / Novelty only
  & F-AUC$\uparrow$
  & 0.8949
  & \textbf{0.8987} \\

Clean / Uncertainty only
  & F-AUC$\uparrow$
  & 0.8953
  & \textbf{0.8987} \\

Palm hole / w/o miss-ray interp.
  & R-AUC$\uparrow$
  & 0.5498
  & $\mathbf{0.5715}$ \\

Pose jitter ($6^\circ$/3\,mm) / w/o RF integration
  & Corr.$\uparrow$
  & 0.2724
  & $\mathbf{0.5028}$ \\
\bottomrule
\end{tabular}

\vspace{1pt}
{\scriptsize
Full = Ray-GPIS. Metrics are failure-specific and not cross-row comparable.
}

\vspace{-0.12cm}
\end{table}

\textbf{Q3 (AC, R2, R1): Is the improved mesh useful downstream?}
Yes. As shown in Table~\ref{tab:downstream}, meshes reconstructed by Ray-GPIS
support placement and subsequent regrasp more reliably than single-view and
non-active reconstruction baselines. Planning uses only the reconstructed
mesh, while GT is used only to evaluate perturbed execution; the GT-mesh
oracle confirms that the task is reachable, and prior real AURORA meshes show
the same downstream applicability.


\begin{table}[t]
\centering
\caption{Downstream success on diverse objects (30 trials/object/method).
}
\label{tab:downstream}
\vspace{-0.22cm}
\resizebox{\linewidth}{!}{%
\begin{tabular}{lccc}
\toprule
Reconstruction source & Placement$\uparrow$ & Regrasp$\uparrow$ & Joint$\uparrow$ \\
\midrule
Single RGB-D & 39.7 & 20.0 & 6.7 \\
SPAR3D (oracle-align) & 30.0 & 15.0 & 5.0 \\
TRELLIS.2 (oracle-align) & 19.7 & 7.7 & 0.0 \\
\midrule
Fixed schedule & 50.0 & 44.0 & 24.0 \\
Adapted PB-NBV & 43.3 & 40.0 & 23.0 \\
Adapted ActNeRF & 56.0 & 72.0 & 41.7 \\
\textbf{Full Ray-GPIS} & \textbf{57.7} & \textbf{77.7} & \textbf{45.0} \\
\midrule
GT mesh (oracle) & 100.0 & 99.0 & 99.0 \\
\bottomrule
\end{tabular}}
\vspace{-0.25cm}
\end{table}

\section*{2. Ablations, Robustness, Diversity \& Runtime}

\textbf{Q4 (R1, R3): Can the key Ray-GPIS components be ablated?}
Yes. Full Ray-GPIS slightly improves overall reconstruction AUC over both
Novelty Only and Uncertainty Only. Its advantage is more
pronounced under targeted failures, miss-ray interpolation improves recovery
of surfaces missing under palm occlusion, while receptive-field integration
maintains a more stable score map under pose jitter
(Table~\ref{tab:robust}(b)).

\textbf{Q5 (R1, AC): How robust is AURORA to 6D pose-tracking errors?}
Yes; on real RGB-D sequences, AURORA retains 96.3\% and 90.3\% of its
zero-error F@5 under $5^\circ$/5\,mm and $10^\circ$/10\,mm pose perturbations,
respectively, showing graceful degradation as 6D pose error increases.

\textbf{Q6 (R3, AC): Does the method generalize to more challenging
irregular/concave objects?} Yes. We add a concave Bowl and a Thin-Irregular
object (eight objects total), each evaluated from 15 poses. Their F-AUC/final
F@5 are $0.6402{\pm}0.0534$/$0.8107{\pm}0.0695$ and
$0.7046{\pm}0.0306$/$0.9274{\pm}0.0219$, respectively.

\textbf{Q7 (R3, R1, AC): What is the per-module runtime, and why replan
only every 6\,s?} The 6\,s interval is required by the hand controller: it needs
this time to respond, complete the tactile rotation, and stabilize the grasp;
replanning earlier would change the command before the hand can reliably
execute it. Segmentation and end-to-end perception/BundleTrack take
$14.5{\pm}0.2$ and $242.9{\pm}47.9$\,ms/frame (the latter includes the former),
while fusion, GP update, candidate scoring, and NBV-to-action mapping take
$1,649.6{\pm}194.0$, $261{\pm}24$, $0.072{\pm}0.001$, and $2.42$\,ms/update,
respectively.

\section*{Other Responses}
\textbf{Q8 (R2): Why use in-hand rotation if Ray-GPIS is
controller-agnostic?} We study in-hand rotation because active reconstruction
under hand occlusion and contact-constrained motion remains largely unexplored.
Each target view must be reached without releasing the object using only
feasible tactile rotations; controller-agnostic describes Ray-GPIS's
transferability, not the absence of this setting-specific challenge.


{\footnotesize \bibliographystyle{plainnat}\bibliography{references}}

\end{document}
